\documentclass[12pt]{article}
\usepackage[utf8]{inputenc}
\usepackage[T1]{fontenc}
\usepackage{hyperref}
\title{Mingrelian Language Resources for K-12 Curriculum Development}
\author{Compiled by Claude Sonnet 5.5 (claude-sonnet-5-5) for Andr\'as Kornai}
\date{}
\begin{document}
\maketitle

\section{Introduction}
Mingrelian (Megrelian; ISO 639-3 xmf; Glottocode ming1252) is a Kartvelian language of western Georgia (Samegrelo) with roughly 350,000 speakers. It has traditionally been non-literary, with Georgian used for writing, and writing in Mingrelian remains debated. Most available materials are scholarly. The searches found few classroom textbooks.

\section{Primary Educational Resources}
\begin{itemize}
\item \emph{Nanashi nina} by Giorgi Sichinava (2002), a children's book teaching the alphabet, as listed in the Wikipedia article: \url{https://en.wikipedia.org/wiki/Mingrelian_language}
\end{itemize}
No textbooks or school programmes were found in the results.

\section{Linguistic and Reference Materials}
\begin{itemize}
\item Wikipedia, Mingrelian language: \url{https://en.wikipedia.org/wiki/Mingrelian_language}
\item Wikipedia, Mingrelian grammar: \url{https://en.wikipedia.org/wiki/Mingrelian_grammar}
\item Reference Grammar for Megrelian, Lund University project (corpus-based): \url{https://projekt.ht.lu.se/megrelian/}
\item EURALEX paper on compiling a bilingual Megrelian-English online dictionary (Lobzhanidze, Gersamia, Tsulaia): \url{https://euralex.org/publications/compiling-a-bilingual-megrelian-english-online-dictionary-preserving-endangered-kartvelian-languages/}
\item Lexilogos Mingrelian dictionary page, which links dictionaries such as Kajaia's Megrelian-Georgian dictionary (2004): \url{https://www.lexilogos.com/english/mingrelian_dictionary.htm}
\item Lexilogos, French-language Mingrelian page: \url{https://www.lexilogos.com/mingrelien_dictionnaire.htm}
\item Rusudan Gersamia, Ilia State University faculty page (researcher on Megrelian): \url{https://faculty.iliauni.edu.ge/en/sas_faculty/rusudan-gersamia/}
\end{itemize}

\section{Religious and Community Resources}
None were found in the search results.

\section{Digital and Audio Resources}
The EURALEX paper describes an annotated corpus and online dictionary built in the Samegrelo documentation project. Whether the corpus is publicly accessible should be checked at the project pages above.

\section{Curriculum Development Context}
Mingrelian is spoken widely but has no established school tradition. Orthography is contested, and Georgian is the written language. A curriculum could build on the documentation project's dictionary and corpus and the Lund reference grammar. Teachers and community members would need to settle a written standard.
\end{document}